% ══════════════════════════════════════════════════════════════════
% CHAPITRE 4 — RÉSULTATS ET DISCUSSION
% Plan : Intro · 4.1 → 4.6 · Conclusion
% ══════════════════════════════════════════════════════════════════

\pagechapitre{chap:resultats}{Résultats et Discussion}

\section*{Introduction}
\addcontentsline{toc}{section}{Introduction}

Ce dernier chapitre présente les résultats obtenus en appliquant
le protocole de tests défini à la section~\ref{sec:protocole}.
Les performances de chacun des trois modules d'IA sont d'abord
exposées et analysées, puis les tests fonctionnels du module
Historique sont restitués. Une discussion générale met ensuite
ces résultats en perspective, avant la présentation du
déploiement de l'application.

% ══════════════════════════════════════════════════════════════════
\section{Résultats du module Irrigation}
\label{sec:res_irrigation}
% ══════════════════════════════════════════════════════════════════

Le tableau~\ref{tab:res_irr} présente les performances du Random
Forest sur l'ensemble de test.

\begin{table}[H]
	\centering
	\caption{Performances du module Irrigation (Random Forest)}
	\label{tab:res_irr}
	\begin{tabularx}{\textwidth}{lXr}
		\toprule
		\textbf{Métrique} & \textbf{Description} & \textbf{Valeur} \\
		\midrule
		Accuracy       & Taux de prédictions correctes      & 78,51\,\% \\
		Précision      & Fiabilité des alertes d'irrigation & 68,59\,\% \\
		Rappel         & Besoins d'irrigation détectés      & 77,03\,\% \\
		F1-Score       & Équilibre précision/rappel         & 72,57\,\% \\
		AUC-ROC        & Pouvoir discriminant               & 0,8782 \\
		CV (k=5)       & Score moyen en validation croisée  & 78,87\,\% \\
		\midrule
		Observations   & Jeu NASA POWER (5 villes)          & 16~435 \\
		Arbres         & \texttt{n\_estimators}             & 100 \\
		\bottomrule
	\end{tabularx}
\end{table}

L'AUC-ROC de 0,8782 traduit une bonne capacité à distinguer les
situations nécessitant une irrigation. Le rappel de 77\,\% ---
privilégié par la pondération équilibrée des classes --- limite
les faux négatifs, c'est-à-dire les stress hydriques non
détectés, plus dommageables pour la culture qu'une irrigation
superflue. L'écart minime entre accuracy de test (78,51\,\%) et
score de validation croisée (78,87\,\%) confirme la stabilité du
modèle.

L'analyse de l'importance des variables place en tête l'humidité
de l'air (0,124), le cumul de pluie sur 7 jours (0,117) et le
cumul sur 14 jours (0,114) : la décision d'irriguer est donc
principalement gouvernée par l'état hydrique récent de
l'atmosphère et du sol, ce qui est agronomiquement cohérent.

% ══════════════════════════════════════════════════════════════════
\section{Résultats du module Fertilisation}
\label{sec:res_fertilisation}
% ══════════════════════════════════════════════════════════════════

\begin{table}[H]
	\centering
	\caption{Performances du module Fertilisation (XGBoost)}
	\label{tab:res_fert}
	\begin{tabularx}{\textwidth}{lXr}
		\toprule
		\textbf{Métrique} & \textbf{Description} & \textbf{Valeur} \\
		\midrule
		Accuracy        & Taux de recommandations correctes  & 95,42\,\% \\
		F1-Score macro  & Moyenne des F1 des 5 classes       & 95,47\,\% \\
		CV (k=5)        & Score moyen en validation croisée  & 97,23\,\% \\
		Écart train/val & Contrôle du surapprentissage       & 3,75\,\% \\
		\midrule
		Observations    & Kaggle réel + augmentation          & 1~599 \\
		Classes         & Types d'engrais                     & 5 \\
		\bottomrule
	\end{tabularx}
\end{table}

Les excellentes performances (tableau~\ref{tab:res_fert})
s'expliquent par la structure du problème : chaque engrais
possède une signature NPK distinctive que le modèle apprend
aisément. L'analyse d'importance le confirme : le phosphore
domine (0,439), suivi de l'azote (0,336) et du potassium
(0,050) --- le phosphore est en effet le principal discriminant
du DAP (18-46-0) face à l'urée (46-0-0) et au KCl (0-0-60).
Le faible écart entraînement/validation (3,75\,\%) écarte
l'hypothèse d'un surapprentissage malgré la part de données
augmentées.

% ══════════════════════════════════════════════════════════════════
\section{Résultats du module Détection des maladies}
\label{sec:res_maladies}
% ══════════════════════════════════════════════════════════════════

\begin{table}[H]
	\centering
	\caption{Performances du CNN MobileNetV2 (dataset tchadien)}
	\label{tab:res_cnn}
	\begin{tabularx}{\textwidth}{lXr}
		\toprule
		\textbf{Métrique} & \textbf{Description} & \textbf{Valeur} \\
		\midrule
		Val. accuracy (phase 1) & Tête seule, base gelée      & 65,12\,\% \\
		Val. accuracy (phase 2) & Après \textit{fine-tuning}  & 60,47\,\% \\
		Accuracy (test)         & Ensemble de test final      & 52,27\,\% \\
		F1-Score macro          & Moyenne des 10 classes      & 47,74\,\% \\
		\midrule
		Images                  & Dataset local tchadien      & 287 \\
		Classes                 & 5 cultures $\times$ 2 états & 10 \\
		Modèle TFLite           & Taille après quantification & $<$ 10 Mo \\
		\bottomrule
	\end{tabularx}
\end{table}

L'accuracy de test de 52,27\,\% --- soit plus de cinq fois le
niveau du hasard (10\,\% pour 10 classes) --- doit être
interprétée au regard des contraintes du dataset :
environ 28 images par classe, très en deçà des standards du
domaine, et des prises de vue en conditions réelles (fonds
complexes, éclairages variables). Cette lecture est confortée
par la littérature : le modèle de Mohanty \textit{et al.}
\cite{mohanty2016deep}, qui atteint 99,35\,\% sur les images de
laboratoire de PlantVillage, chute à environ 31\,\% sur des
images de terrain. Notre résultat, obtenu directement sur des
images de terrain tchadiennes, constitue ainsi une base
honorable, que l'enrichissement du dataset devrait porter à
70--80\,\%. Les classes les mieux reconnues sont
\og Arachide saine \fg{}, \og Maïs sain \fg{} et
\og Maïs malade \fg{} ; les confusions se concentrent sur les
céréales à morphologie proche (mil et sorgho), dont les feuilles
sont difficiles à distinguer même pour un \oe il non expert.

% ══════════════════════════════════════════════════════════════════
\section{Résultats des tests du module Historique}
\label{sec:res_historique}
% ══════════════════════════════════════════════════════════════════

Le plan de tests fonctionnels défini à la
section~\ref{sec:protocole} a été exécuté ; le
tableau~\ref{tab:tests_hist} en présente la synthèse.

\begin{table}[H]
	\centering
	\caption{Résultats des tests fonctionnels du module Historique}
	\label{tab:tests_hist}
	\small
	\begin{tabularx}{\textwidth}{lXc}
		\toprule
		\textbf{Cas de test} & \textbf{Résultat attendu} &
		\textbf{Verdict} \\
		\midrule
		Sauvegarde après analyse & Une ligne insérée dans la table du
		module concerné, horodatée & Validé \\
		Consultation unifiée & Les analyses des trois tables
		apparaissent triées par date décroissante & Validé \\
		Statistiques & Compteurs par module, culture la plus analysée,
		zone à risque et taux calculés exacts & Validé \\
		Export CSV & Fichiers téléchargés lisibles dans un tableur,
		encodage UTF-8, toutes colonnes présentes & Validé \\
		Purge sélective & Suppression par module sans affecter les
		autres tables ; purge totale effective & Validé \\
		Persistance & Les données survivent au redémarrage local de
		l'application & Validé \\
		\bottomrule
	\end{tabularx}
\end{table}

Une limite est à signaler : sur l'hébergement cloud gratuit, le
système de fichiers est éphémère --- la base est réinitialisée à
chaque redémarrage du conteneur. L'export CSV prend alors tout
son sens, en permettant à l'utilisateur de conserver localement
ses données ; en local, la persistance est totale.

% ══════════════════════════════════════════════════════════════════
\section{Discussion générale des résultats}
\label{sec:discussion}
% ══════════════════════════════════════════════════════════════════

\begin{table}[H]
	\centering
	\caption{Synthèse des performances des trois modules d'IA}
	\label{tab:synthese}
	\begin{tabularx}{\textwidth}{lXrr}
		\toprule
		\textbf{Module} & \textbf{Algorithme} & \textbf{Accuracy} &
		\textbf{F1-Score} \\
		\midrule
		Irrigation    & Random Forest    & 78,51\,\% & 72,57\,\% \\
		Fertilisation & XGBoost          & 95,42\,\% & 95,47\,\% \\
		Maladies      & CNN MobileNetV2  & 52,27\,\% & 47,74\,\% \\
		\bottomrule
	\end{tabularx}
\end{table}

Trois enseignements se dégagent du tableau~\ref{tab:synthese}.
D'abord, l'adéquation entre algorithme et nature des données est
déterminante : sur données tabulaires abondantes (irrigation) ou
bien structurées (fertilisation), les méthodes d'ensemble
atteignent des niveaux directement exploitables ; sur images de
terrain en faible quantité, l'apprentissage profond reste
pénalisé malgré le transfert d'apprentissage. Ensuite, l'objectif
de fiabilité fixé au chapitre 2 (accuracy $\geq 75\,\%$) est
atteint pour l'irrigation et largement dépassé pour la
fertilisation. Enfin, la principale marge de progression ---
le module maladies --- relève de la donnée plus que du modèle :
c'est l'enrichissement du dataset local qui conditionnera le
saut de performance.

% ══════════════════════════════════════════════════════════════════
\section{Déploiement de l'application}
\label{sec:deploiement}
% ══════════════════════════════════════════════════════════════════

\subsection{Déploiement local}

En environnement local (Ubuntu 24.04, Python 3.12),
l'application se lance par la commande
\texttt{streamlit run app.py} après installation des dépendances
du fichier \texttt{requirements.txt} dans un environnement
virtuel. L'ensemble des modules, y compris le CNN via
\texttt{tflite-runtime}, y est opérationnel.

\subsection{Déploiement sur Streamlit Cloud}

Le code source est hébergé sur GitHub et déployé en continu sur
Streamlit Cloud \cite{agroiatchad2026} : chaque
\texttt{git push} déclenche la reconstruction de l'application.
Deux adaptations ont été nécessaires : le verrouillage de la
version de Python à 3.11 (fichier \texttt{.python-version}),
et le remplacement de TensorFlow --- indisponible pour les
versions récentes de Python de la plateforme --- par la
bibliothèque légère \texttt{tflite-runtime} pour l'inférence du
CNN. L'application est accessible publiquement à l'adresse :

\begin{center}
	\url{https://agro-ia-tchad.streamlit.app}
\end{center}

\section*{Conclusion}
\addcontentsline{toc}{section}{Conclusion}

Les résultats confirment l'atteinte des objectifs : le module
Fertilisation excelle (95,42\,\%), le module Irrigation dépasse
le seuil de fiabilité fixé (78,51\,\%, AUC 0,878), le module
Maladies fournit une base prometteuse sur images de terrain
réelles (52,27\,\%), et le module Historique remplit
intégralement son cahier des charges de traçabilité et d'export.
L'application, déployée et accessible en ligne, matérialise
l'aboutissement opérationnel du projet. La conclusion générale
qui suit en dresse le bilan et trace les perspectives.